\documentclass[10pt,a4paper]{amsart}
\usepackage{amsmath,amssymb,amsthm,graphicx}
\usepackage{verbatim}
\newenvironment{reference}{\comment}{\endcomment}
%\newenvironment{reference}{}{}
\newenvironment{slogan}{\comment}{\endcomment}
\newenvironment{history}{\comment}{\endcomment}

% For commutative diagrams we use Xy-pic
\usepackage[all]{xy}

% We use 2cell for 2-commutative diagrams.
\xyoption{2cell}
% Theorem environments.
%
\theoremstyle{plain}
\newtheorem{theorem}[subsection]{Theorem}
\newtheorem{proposition}[subsection]{Proposition}
\newtheorem{lemma}[subsection]{Lemma}

\theoremstyle{definition}
\newtheorem{definition}[subsection]{Definition}
\newtheorem{example}[subsection]{Example}
\newtheorem{exercise}[subsection]{Exercise}
\newtheorem{situation}[subsection]{Situation}

\theoremstyle{remark}
\newtheorem{remark}[subsection]{Remark}
\newtheorem{remarks}[subsection]{Remarks}

\numberwithin{equation}{subsection}

% Macros
%
\def\lim{\mathop{\mathrm{lim}}\nolimits}
\def\colim{\mathop{\mathrm{colim}}\nolimits}
\def\Spec{\mathop{\mathrm{Spec}}}
\def\Hom{\mathop{\mathrm{Hom}}\nolimits}
\def\Ext{\mathop{\mathrm{Ext}}\nolimits}
\def\SheafHom{\mathop{\mathcal{H}\!\mathit{om}}\nolimits}
\def\SheafExt{\mathop{\mathcal{E}\!\mathit{xt}}\nolimits}
\def\Sch{\mathit{Sch}}
\def\Mor{\mathop{\mathrm{Mor}}\nolimits}
\def\Ob{\mathop{\mathrm{Ob}}\nolimits}
\def\Sh{\mathop{\mathit{Sh}}\nolimits}
\def\NL{\mathop{N\!L}\nolimits}
\def\CH{\mathop{\mathrm{CH}}\nolimits}
\def\proetale{{pro\text{-}\acute{e}tale}}
\def\etale{{\acute{e}tale}}
\def\QCoh{\mathit{QCoh}}
\def\Ker{\mathop{\mathrm{Ker}}}
\def\Im{\mathop{\mathrm{Im}}}
\def\Coker{\mathop{\mathrm{Coker}}}
\def\Coim{\mathop{\mathrm{Coim}}}

% Boxtimes
%
\DeclareMathSymbol{\boxtimes}{\mathbin}{AMSa}{"02}

%
% Macros for moduli stacks/spaces
%
\def\QCohstack{\mathcal{QC}\!\mathit{oh}}
\def\Cohstack{\mathcal{C}\!\mathit{oh}}
\def\Spacesstack{\mathcal{S}\!\mathit{paces}}
\def\Quotfunctor{\mathrm{Quot}}
\def\Hilbfunctor{\mathrm{Hilb}}
\def\Curvesstack{\mathcal{C}\!\mathit{urves}}
\def\Polarizedstack{\mathcal{P}\!\mathit{olarized}}
\def\Complexesstack{\mathcal{C}\!\mathit{omplexes}}
% \Pic is the operator that assigns to X its picard group, usage \Pic(X)
% \Picardstack_{X/B} denotes the Picard stack of X over B
% \Picardfunctor_{X/B} denotes the Picard functor of X over B
\def\Pic{\mathop{\mathrm{Pic}}\nolimits}
\def\Picardstack{\mathcal{P}\!\mathit{ic}}
\def\Picardfunctor{\mathrm{Pic}}
\def\Deformationcategory{\mathcal{D}\!\mathit{ef}}
\usepackage{cleveref}
\crefname{originallemma}{original lemma}{original lemmas}
\Crefname{originallemma}{Original lemma}{Original lemmas}
\makeatletter
\newlabel{duality-lemma-gorenstein-CM}{{0C00}{source}{}{sourcepointer-1}{}}
\newlabel{duality-lemma-gorenstein-CM@cref}{{[originallemma][0][]0C00}{[1][0][]source}{}{}{}}
\makeatother
\crefname{originallemma}{original lemma}{original lemmas}
\Crefname{originallemma}{Original lemma}{Original lemmas}
\makeatletter
\newlabel{duality-lemma-regular-gorenstein}{{0DWG}{source}{}{sourcepointer-2}{}}
\newlabel{duality-lemma-regular-gorenstein@cref}{{[originallemma][0][]0DWG}{[1][0][]source}{}{}{}}
\makeatother
\crefname{originalsection}{original section}{original sections}
\Crefname{originalsection}{Original section}{Original sections}
\makeatletter
\newlabel{duality-section-gorenstein}{{0AWV}{source}{}{sourcepointer-3}{}}
\newlabel{duality-section-gorenstein@cref}{{[originalsection][0][]0AWV}{[1][0][]source}{}{}{}}
\makeatother
\crefname{originalequation}{original equation}{original equations}
\Crefname{originalequation}{Original equation}{Original equations}
\makeatletter
\newlabel{equation-degree-c1}{{0BRY}{source}{}{sourcepointer-4}{}}
\newlabel{equation-degree-c1@cref}{{[originalequation][0][]0BRY}{[1][0][]source}{}{}{}}
\makeatother
\crefname{originalequation}{original equation}{original equations}
\Crefname{originalequation}{Original equation}{Original equations}
\makeatletter
\newlabel{equation-rr}{{0BRZ}{source}{}{sourcepointer-5}{}}
\newlabel{equation-rr@cref}{{[originalequation][0][]0BRZ}{[1][0][]source}{}{}{}}
\makeatother
\crefname{originalreference}{original reference}{original references}
\Crefname{originalreference}{Original reference}{Original references}
\makeatletter
\newlabel{lemma-duality-dim-1}{{0BS2}{source}{}{sourcepointer-6}{}}
\newlabel{lemma-duality-dim-1@cref}{{[originalreference][0][]0BS2}{[1][0][]source}{}{}{}}
\makeatother
\crefname{originalreference}{original reference}{original references}
\Crefname{originalreference}{Original reference}{Original references}
\makeatletter
\newlabel{lemma-duality-dim-1-CM}{{0BS3}{source}{}{sourcepointer-7}{}}
\newlabel{lemma-duality-dim-1-CM@cref}{{[originalreference][0][]0BS3}{[1][0][]source}{}{}{}}
\makeatother
\crefname{originalreference}{original reference}{original references}
\Crefname{originalreference}{Original reference}{Original references}
\makeatletter
\newlabel{lemma-euler}{{0BS5}{source}{}{sourcepointer-8}{}}
\newlabel{lemma-euler@cref}{{[originalreference][0][]0BS5}{[1][0][]source}{}{}{}}
\makeatother
\crefname{originalreference}{original reference}{original references}
\Crefname{originalreference}{Original reference}{Original references}
\makeatletter
\newlabel{remark-rework-duality-locally-free}{{0BS4}{source}{}{sourcepointer-9}{}}
\newlabel{remark-rework-duality-locally-free@cref}{{[originalreference][0][]0BS4}{[1][0][]source}{}{}{}}
\makeatother
\title{Riemann-Roch and duality for locally free sheaves on a Gorenstein curve}
\author{The Stacks Project Authors}
\date{stacks-project:curves.tex:local-snapshot-20261012; GNU FDL 1.2 or later}
\begin{document}
\maketitle
\noindent\textit{Editorial note.} All delivered mathematical text is the author's own, copied from the pinned Stacks Project chapter file curves.tex; no statement or proof step was rewritten, reordered or removed.\par\medskip
\noindent\textit{Editorial note.} The counted proof is the author's complete argument: the invertibility step, the self-applied Riemann-Roch computation of deg(omega\_X), the substituted chi formula, and the duality symmetry are all retained verbatim.\par\medskip
\noindent\textit{Editorial note.} The label lemma-rr names the counted statement, and the proof environment immediately following it in the source is delivered verbatim as its proof.\par\medskip
\noindent\textit{Editorial note.} The counted statement is carried by the author's own lemma environment in the chapter source.\par\medskip
\noindent\textit{Allowed standard prerequisites (editorial).}\begin{itemize}
\item Gorenstein implies CM and the dualizing module. Duality for Schemes, Lemma duality-lemma-gorenstein-CM, Section duality-section-gorenstein; Same chapter, Lemma lemma-duality-dim-1-CM. Makes omega\_X a dualizing module and, by the Gorenstein property, an invertible sheaf.
\item The Riemann-Roch relation and Euler characteristics. Same chapter, Equation equation-rr, Equation equation-degree-c1 and Lemma lemma-euler. Supplies chi(E) in terms of the Chern class and chi(O\_X), the identity chi(omega\_X) = - deg(omega\_X), and the base Euler characteristic computation.
\item Duality for locally free sheaves. Same chapter, Remark remark-rework-duality-locally-free. Gives the symmetry dim H\textasciicircum{}i(E) = dim H\textasciicircum{}\{1-i\}(E\textasciicircum{}vee (+) omega\_X) for locally free E.
\end{itemize}\medskip
\ifdefined\hypertarget\hypertarget{sourcepointer-1}{}\fi
\noindent\textit{Source reference.} Original lemma 0C00: sources/stacks-chow/tags.txt, source lines 15070--15070. Prerequisite pointer invoked at proof line(s) 1243; it names a standard fact of a sibling chapter that the counted argument uses, and every step of the counted proof built on it is retained verbatim in the delivered body.\par\medskip
\ifdefined\hypertarget\hypertarget{sourcepointer-2}{}\fi
\noindent\textit{Source reference.} Original lemma 0DWG: sources/stacks-chow/tags.txt, source lines 17396--17396. Prerequisite pointer invoked at proof line(s) 1239; it names a standard fact of a sibling chapter that the counted argument uses, and every step of the counted proof built on it is retained verbatim in the delivered body.\par\medskip
\ifdefined\hypertarget\hypertarget{sourcepointer-3}{}\fi
\noindent\textit{Source reference.} Original section 0AWV: sources/stacks-chow/tags.txt, source lines 13733--13733. Prerequisite pointer invoked at proof line(s) 1248; it names a standard fact of a sibling chapter that the counted argument uses, and every step of the counted proof built on it is retained verbatim in the delivered body.\par\medskip
\ifdefined\hypertarget\hypertarget{sourcepointer-4}{}\fi
\noindent\textit{Source reference.} Original equation 0BRY: sources/stacks-curves/pinned-original-source.tex, source lines 1127--1127. Prerequisite pointer invoked at proof line(s) 1258; it names a standard fact of the same chapter that the counted argument uses, and every step of the counted proof built on it is retained verbatim in the delivered body.\par\medskip
\ifdefined\hypertarget\hypertarget{sourcepointer-5}{}\fi
\noindent\textit{Source reference.} Original equation 0BRZ: sources/stacks-curves/pinned-original-source.tex, source lines 1134--1134. Prerequisite pointer invoked at proof line(s) 1249; it names a standard fact of the same chapter that the counted argument uses, and every step of the counted proof built on it is retained verbatim in the delivered body.\par\medskip
\ifdefined\hypertarget\hypertarget{sourcepointer-6}{}\fi
\noindent\textit{Source reference.} Original reference 0BS2: sources/stacks-curves/pinned-original-source.tex, source lines 796--796. Prerequisite pointer invoked at proof line(s) 1221; it names a standard fact of the same chapter that the counted argument uses, and every step of the counted proof built on it is retained verbatim in the delivered body.\par\medskip
\ifdefined\hypertarget\hypertarget{sourcepointer-7}{}\fi
\noindent\textit{Source reference.} Original reference 0BS3: sources/stacks-curves/pinned-original-source.tex, source lines 864--864. Prerequisite pointer invoked at proof line(s) 1245; it names a standard fact of the same chapter that the counted argument uses, and every step of the counted proof built on it is retained verbatim in the delivered body.\par\medskip
\ifdefined\hypertarget\hypertarget{sourcepointer-8}{}\fi
\noindent\textit{Source reference.} Original reference 0BS5: sources/stacks-curves/pinned-original-source.tex, source lines 1174--1174. Prerequisite pointer invoked at proof line(s) 1254; it names a standard fact of the same chapter that the counted argument uses, and every step of the counted proof built on it is retained verbatim in the delivered body.\par\medskip
\ifdefined\hypertarget\hypertarget{sourcepointer-9}{}\fi
\noindent\textit{Source reference.} Original reference 0BS4: sources/stacks-curves/pinned-original-source.tex, source lines 893--893. Prerequisite pointer invoked at proof line(s) 1261; it names a standard fact of the same chapter that the counted argument uses, and every step of the counted proof built on it is retained verbatim in the delivered body.\par\medskip
\begin{lemma}[Riemann-Roch]
\label{lemma-rr}
Let $X$ be a proper scheme over a field $k$ which is Gorenstein and
equidimensional of dimension $1$. Let $\omega_X$ be as in
Lemma \ref{lemma-duality-dim-1}. Then
\begin{enumerate}
\item $\omega_X$ is an invertible $\mathcal{O}_X$-module,
\item $\deg(\omega_X) = -2\chi(X, \mathcal{O}_X)$,
\item for a locally free $\mathcal{O}_X$-module $\mathcal{E}$
of constant rank we have
$$
\chi(X, \mathcal{E}) = \deg(\mathcal{E}) -
\textstyle{\frac{1}{2}} \text{rank}(\mathcal{E}) \deg(\omega_X)
$$
and $\dim_k(H^i(X, \mathcal{E})) =
\dim_k(H^{1 - i}(X, \mathcal{E}^\vee \otimes_{\mathcal{O}_X} \omega_X))$
for all $i \in \mathbf{Z}$.
\end{enumerate}
\end{lemma}

\begin{proof}
Recall that Gorenstein schemes are Cohen-Macaulay
(Duality for Schemes, Lemma \ref{duality-lemma-gorenstein-CM})
and hence $\omega_X$ is a dualizing module on $X$, see
Lemma \ref{lemma-duality-dim-1-CM}.
It follows more or less from the definition of the Gorenstein property
that the dualizing sheaf is invertible, see
Duality for Schemes, Section \ref{duality-section-gorenstein}.
By (\ref{equation-rr}) applied to $\omega_X$ we have
$$
\chi(X, \omega_X) = 
\deg(c_1(\omega_X) \cap [X]_1) + \chi(X, \mathcal{O}_X)
$$
Combined with Lemma \ref{lemma-euler} this gives
$$
2\chi(X, \mathcal{O}_X) = - \deg(c_1(\omega_X) \cap [X]_1) = - \deg(\omega_X)
$$
the second equality by (\ref{equation-degree-c1}). Putting this back into
(\ref{equation-rr}) for $\mathcal{E}$ gives the displayed formula of the lemma.
The symmetry in dimensions is a consequence of duality for $X$, see
Remark \ref{remark-rework-duality-locally-free}.
\end{proof}

\end{document}
